\documentclass{article}

\input{packages}

\title{Nikiforov_lab4}
\author{ksenerem03 }
\date{October 2024}

\begin{document}

\input{title}
\newpage

\section{Цель работы}
Необходимо разработать программное средство, которое соединяется с сервером JIRA через REST API и строит аналитические отчеты. \\
Аналитический отчет состоит из: 
\begin{itemize}
    \item гистограмма, показывающая время, которое задача провела в открытом состоянии;
    \item диаграммы, которые показывают распределение времени по состояниям задачи;
    \item график, который показывает количество заведенных и закрытых задач в день;
    \item график показывающий общее количество задач для пользователя, в которых он указан как исполнитель и репортер;
    \item гистограмма, отражающая время, которое потратил пользователь на ее выполнение;
    \item график, выражающий количество задач по степени серьезности.
\end{itemize}

\newpage

\section{Аналитические отчеты}
Программное средство было реализовано на языке программирования Python 3.8\\
Код программы расположен в системе контроля версий Git: \\
\href{https://github.com/KsenErem/lab_4_task_analytics_software/tree/main}{https://github.com/KsenErem/lab_4_task_analytics_software/tree/main}

Для того чтобы получить отчет, необходимо нажать на соответствующую кнопку в интерфейсе, показанном на рис. (\ref{ris:интерфейс})

\begin{figure}[H]
\center{\includegraphics[scale=0.4]{img/интерфейс.jpg}}
\caption{Гистограмма времени нахождения задач в открытом состоянии}
\label{ris:интерфейс}
\end{figure}

\subsection{\normalsize Гистограмма времени нахождения задач в открытом состоянии}
Ниже на рис.(\ref{ris:Гистограмма_времени_нахождения_задач_в_открытом_состоянии}) представлена гистограмма, показывающая время от момента открытия задачи до момента закрытия. По оси абсцисс расположено время в днях. По оси ординат расположено суммарное количество задач, которое было в открытом виде соответствующее время. Рассматриваем только закрытые задачи по проекту. 
\begin{figure}[H]
\center{\includegraphics[scale=0.4]{img/Гистограмма времени нахождения задач в открытом состоянии.png}}
\caption{Гистограмма времени нахождения задач в открытом состоянии}
\label{ris:Гистограмма_времени_нахождения_задач_в_открытом_состоянии}
\end{figure}

\subsection{\normalsize Гистограмма времени распределения задач в зависимости от состояния задачи}
Ниже на рис.(\ref{ris:Распределение времени задач в состоянии Open} - \ref{ris:Распределение времени задач в состоянии Resolved}) представлены гистограммы, показывающие распределение времени по состояниям задачи. По оси абсцисс расположено время в днях. По оси ординат расположено суммарное количество задач, которое было открыто соответствующее время. Рассматриваем только закрытые задачи по проекту. 

\begin{figure}[H]
\center{\includegraphics[scale=0.4]{img/Распределение времени задач в состоянии Open.png}}
\caption{Распределение времени задач в состоянии Open}
\label{ris:Распределение времени задач в состоянии Open}
\end{figure}

\begin{figure}[H]
\center{\includegraphics[scale=0.4]{img/Распределение времени задач в состоянии Closed.png}}
\caption{Распределение времени задач в состоянии Closed}
\label{ris:Распределение времени задач в состоянии Closed}
\end{figure}

\begin{figure}[H]
\center{\includegraphics[scale=0.4]{img/Распределение времени задач в состоянии In Progress.png}}
\caption{Распределение времени задач в состоянии In Progress}
\label{ris:Распределение времени задач в состоянии In Progress}
\end{figure}

\begin{figure}[H]
\center{\includegraphics[scale=0.3]{img/Распределение времени задач в состоянии Patch Available.png}}
\caption{Распределение времени задач в состоянии Patch Available}
\label{ris:Распределение времени задач в состоянии Patch Available}
\end{figure}

\begin{figure}[H]
\center{\includegraphics[scale=0.4]{img/Распределение времени задач в состоянии Reopened.png}}
\caption{Распределение времени задач в состоянии Reopened}
\label{ris:Распределение времени задач в состоянии Reopened}
\end{figure}

\begin{figure}[H]
\center{\includegraphics[scale=0.3]{img/Распределение времени задач в состоянии Resolved.png}}
\caption{Распределение времени задач в состоянии Resolved}
\label{ris:Распределение времени задач в состоянии Resolved}
\end{figure}

\subsection{\normalsize График, показывающий количество заведенных и закрытых задач в день}
В качестве временного диапазона рассматривались 2 предыдущих месяца. На оси абсцисс отмечены даты. На оси ординат расположено суммарное количество задач. 
На рис. (\ref{ris:Количество открытых и закрытых задач в проекте KAFKA (за последние 2 месяца)}) изображено отдельно количество закрытых и открытых задач, потому что на рис. (\ref{ris:Количество открытых и закрытых задач в проекте KAFKA (за последние 2 месяца) и накопительный итог}) этот момент плохо виден, но можно наглядно увидеть накопительный итог по задачам. 

\begin{figure}[H]
\center{\includegraphics[scale=0.4]{img/Количество открытых и закрытых задач в проекте KAFKA (за последние 2 месяца).png}}
\caption{Количество открытых и закрытых задач в проекте KAFKA (за последние 2 месяца) и накопительный итог по задачам}
\label{ris:Количество открытых и закрытых задач в проекте KAFKA (за последние 2 месяца) и накопительный итог}
\end{figure}

\begin{figure}[H]
\center{\includegraphics[scale=0.4]{img/Количество открытых и закрытых задач в проекте KAFKA (за последние 2 месяца)_1.png}}
\caption{Количество открытых и закрытых задач в проекте KAFKA (за последние 2 месяца)}
\label{ris:Количество открытых и закрытых задач в проекте KAFKA (за последние 2 месяца)}
\end{figure}

\subsection{\normalsize График, показывающий общее количество задач для пользователей, в которых он указан как исполнитель и репортер}
На рис. (\ref{ris:Топ-30 пользователей с максимальным количеством задач в проекте KAFKA.png}) изображена горизонтальная гистограмма, показывающая топ-30 пользователей с максимальным количеством задач. Но оси абсцисс отмечено количество задач. На оси ординат - имя пользователя. 
\begin{figure}[H]
\center{\includegraphics[scale=0.3]{img/Топ-30 пользователей с максимальным количеством задач в проекте KAFKA.png}}
\caption{Топ-30 пользователей с максимальным количеством задач в проекте KAFKA}
\label{ris:Топ-30 пользователей с максимальным количеством задач в проекте KAFKA.png}
\end{figure}

\subsection{\normalsize Гистограмма, отражающая время, которое потратил пользователь на ее выполнение}
На рис.(\ref{ris:Время потраченное на задачу для одного пользователя}) изображена гистограмма, показывающая количество часов, которое потратил пользователь на решение задачи. По оси абсцисс расположено время в часах. По оси ординат - количество задач, соответствующих данному времени.

\begin{figure}[H]
\center{\includegraphics[scale=0.5]{img/Время потраченное на задачу для одного пользователя.png}}
\caption{Зависимость количества задач от потраченного времени на решение этих задач}
\label{ris:Время потраченное на задачу для одного пользователя}
\end{figure}

\subsection{\normalsize График, выражающий количество задач по степени серьезности}
На рис.(\ref{ris:График количества задач по степени серьезности}) изображен график, показывающий количество задач в зависимости от степени серьезности. 

\begin{figure}[H]
\center{\includegraphics[scale=0.3]{img/График количества задач по степени серьезности.png}}
\caption{Зависимость количества задач от потраченного времени на решение этих задач}
\label{ris:График количества задач по степени серьезности}
\end{figure}

Также, было создано 5 модульных автоматических тестов и дистрибутив для установки программного обеспечения на Windows. 

\newpage
\section{Заключение}
В ходе выполнения данной лабораторной работы была разработана программная система для аналитики задач, взаимодействующая с Jira через REST API. Основными результатами работы стали построение аналитических графиков и гистограмм, которые отражают временные и количественные характеристики выполнения задач в проекте.

\end{document}
